
% PIGEAN VERSION 11182025
% Template Version: Version 3 October 2023
% See section 11 of the User Manual for version history
%
%%%%%%%%%%%%%%%%%%%%%%%%%%%%%%%%%%%%%%%%%%%%%%%%%%%%%%%%%%%%%%%%%%%%%%
%%                                                                 %%
%% Please do not use \input{...} to include other tex files.       %%
%% Submit your LaTeX manuscript as one .tex document.              %%
%%                                                                 %%
%% All additional figures and files should be attached             %%
%% separately and not embedded in the \TeX\ document itself.       %%
%%                                                                 %%
%%%%%%%%%%%%%%%%%%%%%%%%%%%%%%%%%%%%%%%%%%%%%%%%%%%%%%%%%%%%%%%%%%%%%

%%\documentclass[referee,sn-basic]{sn-jnl}% referee option is meant for double line spacing

%%=======================================================%%
%% to print line numbers in the margin use lineno option %%
%%=======================================================%%

%%\documentclass[lineno,sn-basic]{sn-jnl}% Basic Springer Nature Reference Style/Chemistry Reference Style

%%======================================================%%
%% to compile with pdflatex/xelatex use pdflatex option %%
%%======================================================%%

%%\documentclass[pdflatex,sn-basic]{sn-jnl}% Basic Springer Nature Reference Style/Chemistry Reference Style


%%Note: the following reference styles support Namedate and Numbered referencing. By default the style follows the most common style. To switch between the options you can add or remove “Numbered” in the optional parenthesis. 
%%The option is available for: sn-basic.bst, sn-vancouver.bst, sn-chicago.bst%  
 
\documentclass[sn-nature,lineno]{sn-jnl}% Style for submissions to Nature Portfolio journals
%%\documentclass[sn-basic]{sn-jnl}% Basic Springer Nature Reference Style/Chemistry Reference Style
%%\documentclass[sn-mathphys-num]{sn-jnl}% Math and Physical Sciences Numbered Reference Style 
%%\documentclass[sn-mathphys-ay]{sn-jnl}% Math and Physical Sciences Author Year Reference Style
%%\documentclass[sn-aps]{sn-jnl}% American Physical Society (APS) Reference Style
%%\documentclass[sn-vancouver,Numbered]{sn-jnl}% Vancouver Reference Style
%%\documentclass[sn-apa]{sn-jnl}% APA Reference Style 
%%\documentclass[sn-chicago]{sn-jnl}% Chicago-based Humanities Reference Style

%%%% Standard Packages
%%<additional latex packages if required can be included here>

% Changing geoemtry to look better (imo)
\geometry{a4paper, margin=1in}

% Change font to arial 
\usepackage{helvet}
\renewcommand{\familydefault}{\sfdefault}

\usepackage{graphicx}%
\usepackage{multirow}%
\usepackage{amsmath,amssymb,amsfonts}%
\usepackage{amsthm}%
\usepackage{mathrsfs}%
\usepackage[title]{appendix}%
\usepackage{xcolor}%
\usepackage{textcomp}%
\usepackage{manyfoot}%
\usepackage{booktabs}%
\usepackage{longtable}%
\usepackage{algorithm}%
\usepackage{algorithmicx}%
\usepackage{algpseudocode}%
\usepackage{listings}%
\usepackage{subcaption}
\usepackage{float}
\usepackage{placeins}
\usepackage{comment}
\usepackage{tikz}
\usetikzlibrary{positioning,arrows}
% \usepackage{luacode} % Removed: requires LuaTeX, not compatible with pdflatex
\usepackage{pdfpages}
\usepackage{hyperref}
%%%%

%%%%%=============================================================================%%%%
%%%%  Remarks: This template is provided to aid authors with the preparation
%%%%  of original research articles intended for submission to journals published 
%%%%  by Springer Nature. The guidance has been prepared in partnership with 
%%%%  production teams to conform to Springer Nature technical requirements. 
%%%%  Editorial and presentation requirements differ among journal portfolios and 
%%%%  research disciplines. You may find sections in this template are irrelevant 
%%%%  to your work and are empowered to omit any such section if allowed by the 
%%%%  journal you intend to submit to. The submission guidelines and policies 
%%%%  of the journal take precedence. A detailed User Manual is available in the 
%%%%  template package for technical guidance.
%%%%%=============================================================================%%%%

%% as per the requirement new theorem styles can be included as shown below
\theoremstyle{thmstyleone}%
\newtheorem{theorem}{Theorem}%  meant for continuous numbers
%%\newtheorem{theorem}{Theorem}[section]% meant for sectionwise numbers
%% optional argument [theorem] produces theorem numbering sequence instead of independent numbers for Proposition
\newtheorem{proposition}[theorem]{Proposition}% 
%%\newtheorem{proposition}{Proposition}% to get separate numbers for theorem and proposition etc.

\theoremstyle{thmstyletwo}%
\newtheorem{example}{Example}%
\newtheorem{remark}{Remark}%

\theoremstyle{thmstylethree}%
\newtheorem{definition}{Definition}%

%% New commands
% Probablity Dist function
\newcommand{\prob}{\text{Pr}}
% Variable lookup function
\newcommand{\getvar}[1]{\directlua{lookup("#1")}}

\raggedbottom
%%\unnumbered% uncomment this for unnumbered level heads

\begin{document}

\title[Article Title]{Indirect Genetic Support: A Probabilistic Framework for Distinguishing Disease-Relevant from Disease-Associated Genes}

% Alternative titles:
% - A Probabilistic Evidence Map of Gene-Trait Relevance Across the Phenome
% - Calibrating Non-Genetic Evidence with Genetics to Map Gene-Trait Relevance
% - PIGEAN: Using Genetic Associations to Calibrate Functional Annotations Across the Phenome


%%=============================================================%%
%% GivenName	-> \fnm{Joergen W.}
%% Particle	-> \spfx{van der} -> surname prefix
%% FamilyName	-> \sur{Ploeg}
%% Suffix	-> \sfx{IV}
%% \author*[1,2]{\fnm{Joergen W.} \spfx{van der} \sur{Ploeg} 
%%  \sfx{IV}}\email{iauthor@gmail.com}
%%=============================================================%%

\author*[1,2]{\fnm{First} \sur{Author}}\email{iauthor@gmail.com}

\author[2,3]{\fnm{Second} \sur{Author}}\email{iiauthor@gmail.com}
\equalcont{These authors contributed equally to this work.}

\author[1,2]{\fnm{Third} \sur{Author}}\email{iiiauthor@gmail.com}
\equalcont{These authors contributed equally to this work.}

\affil*[1]{\orgdiv{Department}, \orgname{Organization}, \orgaddress{\street{Street}, \city{City}, \postcode{100190}, \state{State}, \country{Country}}}

\affil[2]{\orgdiv{Department}, \orgname{Organization}, \orgaddress{\street{Street}, \city{City}, \postcode{10587}, \state{State}, \country{Country}}}

\affil[3]{\orgdiv{Department}, \orgname{Organization}, \orgaddress{\street{Street}, \city{City}, \postcode{610101}, \state{State}, \country{Country}}}

%%==================================%%
%% Sample for unstructured abstract %%
%%==================================%%

\abstract{
Drug targets with human genetic support are more than twice as likely to succeed in clinical trials, yet genetic associations directly implicate only a fraction of disease-relevant genes. A gene may also be disease-relevant through \textit{indirect genetic support}: the extent to which it shares functional annotations---such as pathway membership, mouse knockout phenotypes, or expression patterns---with other genes that are directly associated with a trait. Here we introduce PIGEAN, a Bayesian framework that uses genetic associations as an anchor to calibrate {{NUM_TOTAL_GENE_ANNOTATIONS:,}} functional annotations and quantify both direct and indirect genetic support for any gene-trait pair. Applied to {{NUM_TOTAL_TRAITS_MOUSE_MSIGDB:,}} traits, PIGEAN produces {{NUM_GENETRAIT_REL_PP50:,}} calibrated gene-trait relationships at posterior probability above 50\%, expanding the genetically supported landscape by an order of magnitude. We validate these predictions on three independent axes: indirect support predicts held-out genetic associations in leave-one-chromosome-out and temporal analyses ($p = {{TEMPORAL_DELTA_R2_BETA_P}}$), prioritizes drug targets at clinical success rates exceeding those of the Open Targets Platform and comparable to Open Targets Genetics L2G scores (relative success = {{MINIKEL_OTG_RS:.2f}}) while identifying {{NUM_PIGEAN_PRIORITIZED_TARGETS_MOUSE_MSIGDB_AT_RS_eq_MINIKEL_OTG_INDIRECT:,}} targets compared to {{NUM_MINIKEL_OTG_PRIORITIZED_TARGETS:,}} from L2G alone, and enriches for causal genes in rare disease patients. Across the phenome, the resulting resource classifies {{NUM_UNDERSTOOD_MECHANISMS:,}} gene-trait pairs as understood mechanisms with convergent direct and indirect evidence, flags {{NUM_NOVEL_DISCOVERIES:,}} as novel discoveries lacking pathway annotation, and generates mechanistic hypotheses for individual genes---for example, decomposing the \textit{HLA-DRB1} effect in multiple sclerosis into a three-stage immune cascade that explains the known sex paradox in disease progression. These results establish indirect genetic support as a quantitative and scalable measure of gene-trait relevance distinct from genetic association, and provide a dense probabilistic map of gene-trait relationships across the human phenome.
}

%%================================%%
%% Sample for structured abstract %%
%%================================%%

% \abstract{\textbf{Purpose:} The abstract serves both as a general introduction to the topic and as a brief, non-technical summary of the main results and their implications. The abstract must not include subheadings (unless expressly permitted in the journal's Instructions to Authors), equations or citations. As a guide the abstract should not exceed 200 words. Most journals do not set a hard limit however authors are advised to check the author instructions for the journal they are submitting to.
% 
% \textbf{Methods:} The abstract serves both as a general introduction to the topic and as a brief, non-technical summary of the main results and their implications. The abstract must not include subheadings (unless expressly permitted in the journal's Instructions to Authors), equations or citations. As a guide the abstract should not exceed 200 words. Most journals do not set a hard limit however authors are advised to check the author instructions for the journal they are submitting to.
% 
% \textbf{Results:} The abstract serves both as a general introduction to the topic and as a brief, non-technical summary of the main results and their implications. The abstract must not include subheadings (unless expressly permitted in the journal's Instructions to Authors), equations or citations. As a guide the abstract should not exceed 200 words. Most journals do not set a hard limit however authors are advised to check the author instructions for the journal they are submitting to.
% 
% \textbf{Conclusion:} The abstract serves both as a general introduction to the topic and as a brief, non-technical summary of the main results and their implications. The abstract must not include subheadings (unless expressly permitted in the journal's Instructions to Authors), equations or citations. As a guide the abstract should not exceed 200 words. Most journals do not set a hard limit however authors are advised to check the author instructions for the journal they are submitting to.}

\keywords{keyword1, Keyword2, Keyword3, Keyword4}

%%\pacs[JEL Classification]{D8, H51}

%%\pacs[MSC Classification]{35A01, 65L10, 65L12, 65L20, 65L70}

\maketitle

\newpage
\tableofcontents
\newpage

\input{sections/indirect-support/03302026/introduction_03302026}

\section{Results}

% Results story (updated 03/04/2026 - refactored for method-focused paper):
% 1. Framework overview - direct vs indirect support, the PIGEAN model
% 2. Validation: indirect predicts direct (LOCO, temporal, rare disease)
% 3. Context-dependent relevance: which annotations for which contexts
% 4. Drug target prioritization: RS analysis
% 5. Rare disease gene enrichment: clinical validation
% 6. Phenome-wide resource description
%
% NOTE: Discovery cycle / saturation analysis moved to companion paper
% (saturation-letter). See versions/saturation-letter/main_03042026.tex

\input{sections/indirect-support/03302026/methods_overview_03302026}

\input{sections/indirect-support/03302026/validation_genetics_03302026}

\input{sections/indirect-support/03302026/biological_properties_03302026}

\input{sections/indirect-support/03302026/geneset_context_03302026}

\subsection{Indirect Support Prioritizes Drug Targets With Higher Clinical Success}

% STRUCTURE: This section has 4 paragraphs:
% P1: The gap (no non-genetic method predicts drug success) + continuous RS framework + weighted EA
% P2: Resources evaluated and curation
% P3: Main RS curve results (Figure 4a,c) + temporal holdout (Figure 4b)
% P4: Weighted EA by clinical area (Figure 5)
% Discussion-worthy deeper stories (oncology saturation, cardio ion channels, metabolic lipid/glucose)
% are flagged with DISCUSSION tags and written in the discussion section.

\begin{figure}[t!]
    \centering
    {{canva:5}}
    \caption{Drug target validation. (a) Continuous relative success (RS) as a function of the fraction of each resource prioritized ($\alpha$), comparing PIGEAN indirect, direct, and combined support (solid lines) to Open Targets Genetics (OTG) locus-to-gene share (dashed lines) and Open Targets Platform (OTP) combined and indirect support (dotted lines). Shaded bands denote 95\% confidence intervals computed by the delta method. Annotations indicate the number of target--indication (T-I) pairs prioritized at selected thresholds. (b) Temporal holdout: RS for drug targets launched after 2018, scored using PIGEAN with pre-2018 gene set libraries on 67 validation traits. (c) Cross-resource comparison of maximum RS achieved after at least 5\% of each resource is prioritized. For binary resources (Genebass, OMIM, Piccolo), RS is computed from a single supported/unsupported partition. Error bars denote 95\% delta method confidence intervals; numbers above each point indicate the number of T-I pairs in the supported set. Resources are grouped by curation source: Minikel curations (left), MAGMA (center), OTP (center-right), and PIGEAN (right), with score types distinguished by color (legend).}
    \label{fig:validation-pharma}
\end{figure}

% P1: The gap + continuous RS + weighted EA metrics
Minikel et al.\ demonstrated that drug targets with genetic support, defined as a gene--trait association in a genetic resource such as Open Targets Genetics (OTG), are approximately two-fold more likely to progress through clinical trials than unsupported targets \cite{minikel_refining_2024}. This finding has motivated efforts to use human genetic evidence for target selection, but genetic association data alone have limited coverage: even the largest compilation of OTG locus-to-gene (L2G) scores prioritizes only {{MINIKEL_OTG_ALL_N_TARGETS:,}} target--indication (T-I) pairs (Figure~\ref{fig:validation-pharma}c). Non-genetic evidence sources such as biological pathways, model organism phenotypes, and literature co-occurrence could substantially expand the number of evaluable targets. The Open Targets Platform (OTP) Europe PMC literature score is the closest existing non-genetic prioritization metric \cite{buniello_open_2025}, yet it achieves a maximum RS of only {{OTP_INDIRECT_RS:.2f}} (95\% CI: {{OTP_INDIRECT_RS_CI_LOWER:.2f}} to {{OTP_INDIRECT_RS_CI_UPPER:.2f}}) across {{OTP_INDIRECT_N_TARGETS:,}} T-I pairs (Figure~\ref{fig:validation-pharma}c), barely distinguishable from the null expectation of $\text{RS} = 1$. No method has demonstrated that non-genetic evidence can be systematically converted into a continuous score that meaningfully predicts drug target clinical success.

% TODO: CITE_OTP_EUROPEPMC - Add citation for Open Targets Platform Europe PMC score description

PIGEAN's indirect support fills this gap. Because PIGEAN, unlike binary classifiers, assigns continuous scores to gene--trait associations, we developed the \textit{continuous RS curve}, which traces $\text{RS}_{\text{cum}}$ as a function of the fraction $\alpha$ of T-I pairs prioritized (Section~\ref{sec:methods-rs}; Figure~\ref{fig:validation-pharma}a). This curve provides a complete profile of enrichment across all possible score thresholds, avoiding the need to select an arbitrary cutoff. To summarize the curve in a single scalar that balances prioritization quality against resource completeness, we defined the \textit{weighted excess area} (EA) above a baseline RS, computed as the integral of $(\text{RS}_{\text{cum}}(\alpha) - \text{RS}_{\text{baseline}})$ weighted by the number of T-I pairs at each threshold (Section~\ref{sec:methods-rs}). We report two variants: EA above $\text{RS} = 1$ ($\text{EA}|\text{RS}{=}1$), which captures total enrichment for successful targets across all thresholds, and EA above $\text{RS} = \text{Self}$ ($\text{EA}|\text{RS}{=}\text{Self}$), which uses the RS when the entire resource is designated as ``supported'' as the baseline, thereby isolating how much the resource's internal ranking adds beyond simple presence or absence in the resource. Because both metrics are weighted by target count, resources with greater coverage are rewarded: a resource that achieves high RS for many targets will score higher than one with equivalent RS for only a handful.

% P2: Resources evaluated and curation
We evaluated PIGEAN alongside several external resources using the continuous RS framework (Section~\ref{sec:methods-rs}). We matched gene--trait associations to T-I pairs from the Pharmaprojects clinical trial database across {{NUM_TOTAL_TRAITS_MOUSE_MSIGDB}} phenotypes using MeSH-based semantic similarity (threshold $\geq 0.8$). As benchmarks, we included (i) discrete curations by Minikel et al.\ \cite{minikel_refining_2024} from Genebass ($N$ = {{MINIKEL_GENEBASS_N_TARGETS}}), OMIM ($N$ = {{MINIKEL_OMIM_N_TARGETS}}), and Piccolo ($N$ = {{MINIKEL_PICCOLO_N_TARGETS}}), in which a T-I pair is either supported or not; (ii) continuous L2G share scores curated by Minikel et al.\ from OTG across five biobank resources (All OTG, FinnGen, GWAS Catalog, UKBB Neale, Saige); (iii) the Open Targets Platform (OTP) association scores, queried from the BigQuery public dataset filtering to non-drug evidence types (\texttt{datatypeId} $\in$ \{animal\_model, genetic\_association, literature\}), with the \textit{combined} score (harmonic mean of genetic and non-genetic evidence excluding drug-derived sources) and the \textit{indirect} score (Europe PMC literature evidence); (iv) MAGMA gene-level $p$-values \cite{de_leeuw_magma_2015} computed on the same summary statistics used for PIGEAN, restricted to genes with $p < 2.5 \times 10^{-6}$; and (v) PIGEAN posterior probabilities (combined, direct, and indirect) from the Mouse + MSigDB model. For continuous resources in Figure~\ref{fig:validation-pharma}c, the reported RS is the maximum achieved after at least 5\% of the resource is prioritized, ensuring numerical stability; for discrete resources, RS is computed from the single supported/unsupported partition.

% P3: Main RS results
% P3a: Indirect support works -- using genetics as anchor to leverage non-genetic connections
Across all sources, PIGEAN indirect support achieved RS = {{PIGEAN_INDIRECT_RS_ALL:.2f}} (95\% CI: {{PIGEAN_INDIRECT_RS_ALL_CI_LOWER:.2f}} to {{PIGEAN_INDIRECT_RS_ALL_CI_UPPER:.2f}}) while prioritizing {{PIGEAN_INDIRECT_N_TARGETS_ALL:,}} T-I pairs, covering 78\% of OTP indirect's target set ({{OTP_INDIRECT_N_TARGETS:,}}) at 1.8-fold higher RS ({{PIGEAN_INDIRECT_RS_ALL:.2f}} versus {{OTP_INDIRECT_RS:.2f}}; Figure~\ref{fig:validation-pharma}c). PIGEAN combined support (direct plus indirect) achieved RS = {{PIGEAN_COMBINED_RS_ALL:.2f}} (95\% CI: {{PIGEAN_COMBINED_RS_ALL_CI_LOWER:.2f}} to {{PIGEAN_COMBINED_RS_ALL_CI_UPPER:.2f}}; $N$ = {{PIGEAN_COMBINED_N_TARGETS_ALL:,}}), and PIGEAN direct support achieved RS = {{PIGEAN_DIRECT_RS_ALL:.2f}} (95\% CI: {{PIGEAN_DIRECT_RS_ALL_CI_LOWER:.2f}} to {{PIGEAN_DIRECT_RS_ALL_CI_UPPER:.2f}}; $N$ = {{PIGEAN_DIRECT_N_TARGETS_ALL:,}}). Indirect support (RS = {{PIGEAN_INDIRECT_RS_ALL:.2f}}) approached combined support (RS = {{PIGEAN_COMBINED_RS_ALL:.2f}}), indicating that by using genetic associations as an anchor, PIGEAN can leverage shared pathway and annotation structure to prioritize a substantially larger set of targets than genetics alone, at comparable enrichment levels. The indirect score requires direct genetic evidence to calibrate which annotations are trait-relevant (Section~\ref{sec:methods-rs}), but the resulting annotation-based prioritization extends far beyond the directly associated genes, reaching {{PIGEAN_INDIRECT_N_TARGETS_ALL:,}} T-I pairs compared with {{MINIKEL_OTG_ALL_N_TARGETS:,}} from OTG L2G.

% P3b: Comparison to curated resources -- more targets at comparable RS
The continuous RS curve (Figure~\ref{fig:validation-pharma}a) reveals a fundamental tradeoff between enrichment magnitude and target coverage. Among the Minikel curations, the highest RS values were achieved by the smallest, most stringent resources: UKBB Neale L2G share (RS = {{MINIKEL_UKBB_RS:.2f}}; 95\% CI: {{MINIKEL_UKBB_RS_CI_LOWER:.2f}} to {{MINIKEL_UKBB_RS_CI_UPPER:.2f}}; $N$ = {{MINIKEL_UKBB_N_TARGETS}}), OMIM (RS = {{MINIKEL_OMIM_RS:.2f}}; 95\% CI: {{MINIKEL_OMIM_RS_CI_LOWER:.2f}} to {{MINIKEL_OMIM_RS_CI_UPPER:.2f}}; $N$ = {{MINIKEL_OMIM_N_TARGETS}}), and Saige (RS = {{MINIKEL_SAIGE_RS:.2f}}; 95\% CI: {{MINIKEL_SAIGE_RS_CI_LOWER:.2f}} to {{MINIKEL_SAIGE_RS_CI_UPPER:.2f}}; $N$ = {{MINIKEL_SAIGE_N_TARGETS}}). MAGMA achieved RS = {{MAGMA_RS:.2f}} (95\% CI: {{MAGMA_RS_CI_LOWER:.2f}} to {{MAGMA_RS_CI_UPPER:.2f}}) for {{MAGMA_N_TARGETS}} T-I pairs. These resources achieve high RS but for target sets that are one to two orders of magnitude smaller than PIGEAN's: PIGEAN combined support covers 63-fold more T-I pairs than UKBB Neale ({{PIGEAN_COMBINED_N_TARGETS_ALL:,}} versus {{MINIKEL_UKBB_N_TARGETS}}) and 12-fold more than All OTG L2G share ({{PIGEAN_COMBINED_N_TARGETS_ALL:,}} versus {{MINIKEL_OTG_ALL_N_TARGETS:,}}), at only modestly lower RS ({{PIGEAN_COMBINED_RS_ALL:.2f}} versus {{MINIKEL_UKBB_RS:.2f}} and {{MINIKEL_OTG_ALL_RS:.2f}}, respectively). Notably, the OTG L2G curve in Figure~\ref{fig:validation-pharma}a shows relatively flat RS across L2G scores, such that even scores approaching 1.0 do not yield substantially higher RS than scores near 0.1, consistent with the findings of Minikel et al.\ \cite{minikel_refining_2024}. By contrast, PIGEAN's RS curve rises steeply at stringent thresholds (Figure~\ref{fig:validation-pharma}a), indicating that its scoring has greater discriminative power for identifying clinically successful targets.

% P3c: OTP fails to rank targets
OTP's combined and indirect scores, despite covering {{OTP_COMBINED_N_TARGETS:,}} and {{OTP_INDIRECT_N_TARGETS:,}} T-I pairs respectively, achieved RS of only {{OTP_COMBINED_RS:.2f}} (95\% CI: {{OTP_COMBINED_RS_CI_LOWER:.2f}} to {{OTP_COMBINED_RS_CI_UPPER:.2f}}) and {{OTP_INDIRECT_RS:.2f}} (95\% CI: {{OTP_INDIRECT_RS_CI_LOWER:.2f}} to {{OTP_INDIRECT_RS_CI_UPPER:.2f}}), hovering near the null across the full range of thresholds (Figure~\ref{fig:validation-pharma}a, dotted lines). The baseline RS (i.e., when the entire resource is designated as supported) was {{OTP_COMBINED_BASELINE_RS:.2f}} for OTP combined and {{OTP_INDIRECT_BASELINE_RS:.2f}} for OTP indirect, both at or below 1.0, indicating that OTP's full target set is not enriched for clinical success relative to the Pharmaprojects background. This stands in contrast to PIGEAN, where the baseline RS of {{PIGEAN_COMBINED_BASELINE_RS_ALL:.2f}} indicates mild enrichment even before applying any score threshold.

% P3d: Source-matched comparisons
Source-matched comparisons confirm that PIGEAN's advantage is not driven by trait selection. Restricting to GWAS Catalog (GCAT) traits, PIGEAN direct support achieved RS = {{PIGEAN_DIRECT_RS_GCAT:.2f}} ($N$ = {{PIGEAN_DIRECT_N_TARGETS_GCAT:,}}), compared with {{PIGEAN_COMBINED_RS_GCAT:.2f}} for combined ($N$ = {{PIGEAN_COMBINED_N_TARGETS_GCAT:,}}) and {{PIGEAN_INDIRECT_RS_GCAT:.2f}} for indirect ($N$ = {{PIGEAN_INDIRECT_N_TARGETS_GCAT:,}}). For portal traits, indirect and combined support were nearly identical (RS = {{PIGEAN_INDIRECT_RS_PORTAL:.2f}} and {{PIGEAN_COMBINED_RS_PORTAL:.2f}}, respectively; both $N$ = {{PIGEAN_COMBINED_N_TARGETS_PORTAL:,}}), suggesting that for common diseases with full summary statistics, indirect evidence captures nearly all of the signal present in combined support. For Orphanet rare diseases, indirect support achieved the highest RS of any source--trait combination: RS = {{PIGEAN_INDIRECT_RS_ORPHANET:.2f}} (95\% CI: {{PIGEAN_INDIRECT_RS_ORPHANET_CI_LOWER:.2f}} to {{PIGEAN_INDIRECT_RS_ORPHANET_CI_UPPER:.2f}}; $N$ = {{PIGEAN_INDIRECT_N_TARGETS_ORPHANET:,}}), essentially equivalent to combined support (RS = {{PIGEAN_COMBINED_RS_ORPHANET:.2f}}; $N$ = {{PIGEAN_COMBINED_N_TARGETS_ORPHANET:,}}). This result indicates that indirect support is especially effective for rare diseases, where pathway biology is well-characterized but direct genetic evidence may be limited to small numbers of known causal genes.

% P3e: Temporal holdout
To assess whether the observed enrichment is driven by circularity, specifically gene set annotations that encode existing drug target knowledge, we performed a temporal holdout analysis (Section~\ref{sec:methods-rs}; Figure~\ref{fig:validation-pharma}b). PIGEAN was re-run using only gene set libraries curated from pre-2018 data, and evaluation was restricted to T-I pairs for drugs launched after 2018 across 67 validation traits. Under this temporal holdout, indirect support achieved a maximum RS of {{MAX_RS_INDIRECT_TEMPORAL_HOLDOUT:.2f}} (95\% CI: {{MAX_RS_INDIRECT_TEMPORAL_HOLDOUT_CI_LOWER:.2f}} to {{MAX_RS_INDIRECT_TEMPORAL_HOLDOUT_CI_UPPER:.2f}}; $N$ = {{NUM_DRUG_TARGETS_SUPPORTED_INDIRECT_TEMPORAL_HOLDOUT}}) and combined support achieved {{MAX_RS_COMBINED_TEMPORAL_HOLDOUT:.2f}} (95\% CI: {{MAX_RS_COMBINED_TEMPORAL_HOLDOUT_CI_LOWER:.2f}} to {{MAX_RS_COMBINED_TEMPORAL_HOLDOUT_CI_UPPER:.2f}}; $N$ = {{NUM_DRUG_TARGETS_SUPPORTED_COMBINED_TEMPORAL_HOLDOUT}}). At 2\% of the resource, RS stabilized at {{RS_INDIRECT_TEMPORAL_HOLDOUT_STABLE_AT_ALPHA_002:.2f}} for indirect support ($N$ = {{NUM_TARGETS_INDIRECT_TEMPORAL_HOLDOUT_AT_ALPHA_002}}) and {{RS_COMBINED_TEMPORAL_HOLDOUT_STABLE_AT_ALPHA_002:.2f}} for combined support ($N$ = {{NUM_TARGETS_COMBINED_TEMPORAL_HOLDOUT_AT_ALPHA_002}}), confirming that the enrichment is not an artifact of circularity with existing drug target annotations.

\begin{figure}[t!]
    \centering
    {{canva:6}}
    \caption{Weighted excess area (EA) above RS baseline across therapeutic areas. Each cell shows the weighted EA (top) and the number of supported T-I pairs (bottom, in parentheses) for a given method and therapeutic area. Left panel: baseline $\text{RS} = 1$ (comparison against all targets). Right panel: baseline $\text{RS} = \text{Self}$ (comparison against each method's own unsupported targets). Bold values indicate the highest weighted EA within each therapeutic area. Color intensity is scaled per column (therapeutic area) to highlight the best-performing method within each area.}
    \label{fig:weighted-ea}
\end{figure}

% P4: Weighted EA by clinical area
% P4a: Overall weighted EA
To compare resources in a manner that jointly accounts for enrichment magnitude and target coverage, we computed the weighted EA across 17 therapeutic areas (Figure~\ref{fig:weighted-ea}). Across all areas combined, PIGEAN combined support achieved the highest $\text{EA}|\text{RS}{=}1$ ({{WEA_PIGEAN_COMBINED_ALL:,}}; $N$ = {{PIGEAN_COMBINED_N_TARGETS_ALL:,}}), followed by PIGEAN indirect support ({{WEA_PIGEAN_INDIRECT_ALL:,}}; $N$ = {{PIGEAN_INDIRECT_N_TARGETS_ALL:,}}), PIGEAN direct support ({{WEA_PIGEAN_DIRECT_ALL:,}}; $N$ = {{PIGEAN_DIRECT_N_TARGETS_ALL:,}}), OTP combined ({{WEA_OTP_COMBINED_ALL:,}}; $N$ = {{OTP_COMBINED_N_TARGETS:,}}), OTP indirect ({{WEA_OTP_INDIRECT_ALL:,}}; $N$ = {{OTP_INDIRECT_N_TARGETS:,}}), OTG L2G share ({{WEA_OTG_L2G_ALL:,}}; $N$ = {{MINIKEL_OTG_ALL_N_TARGETS:,}}), and MAGMA ({{WEA_MAGMA_ALL}}; $N$ = {{MAGMA_N_TARGETS}}). Notably, PIGEAN combined and indirect support dominated despite covering fewer T-I pairs than OTP ({{PIGEAN_COVERAGE_FRACTION:.1%}} versus {{OTP_COVERAGE_FRACTION:.1%}} of all T-I pairs): because the weighted EA metric rewards coverage, OTP's broader target set should confer an advantage, yet PIGEAN's per-target enrichment signal was strong enough to more than compensate. On the $\text{EA}|\text{RS}{=}\text{Self}$ metric, which isolates ranking ability from baseline target quality, PIGEAN combined support again led ({{WEA_SELF_PIGEAN_COMBINED_ALL:,}}), followed by PIGEAN indirect ({{WEA_SELF_PIGEAN_INDIRECT_ALL:,}}), OTP combined ({{WEA_SELF_OTP_COMBINED_ALL:,}}), OTG L2G share ({{WEA_SELF_OTG_L2G_ALL}}), and MAGMA ({{WEA_SELF_MAGMA_ALL}}). This confirms that PIGEAN's scoring provides genuine ranking information, not merely a better pool of targets, and that indirect support alone captures the majority of this discriminative signal.

% P4b: Areas where PIGEAN dominates
The per-area breakdown reveals that PIGEAN's advantage is concentrated in therapeutic areas where pathway biology provides strong signal (Figure~\ref{fig:weighted-ea}). In oncology, PIGEAN indirect support achieved $\text{EA}|\text{RS}{=}1$ = {{WEA_PIGEAN_INDIRECT_ONCO:,}} ($N$ = {{WEA_PIGEAN_INDIRECT_ONCO_N:,}}), while OTP combined was negative ($\text{EA}|\text{RS}{=}1$ = {{WEA_OTP_COMBINED_ONCO:,}}; $N$ = {{WEA_OTP_COMBINED_ONCO_N:,}}), indicating that OTP's scoring actively anti-predicts oncology drug success. In neurology, PIGEAN combined achieved $\text{EA}|\text{RS}{=}1$ = {{WEA_PIGEAN_COMBINED_NEURO:,}} ($N$ = {{WEA_PIGEAN_COMBINED_NEURO_N:,}}) and PIGEAN indirect achieved {{WEA_PIGEAN_INDIRECT_NEURO:,}}, compared with {{WEA_OTP_COMBINED_NEURO}} for OTP combined ($N$ = {{WEA_OTP_COMBINED_NEURO_N:,}}), a six-fold advantage. In cardiovascular disease, PIGEAN indirect support ($\text{EA}|\text{RS}{=}1$ = {{WEA_PIGEAN_INDIRECT_CARDIO}}; $N$ = {{WEA_PIGEAN_INDIRECT_CARDIO_N:,}}) outperformed PIGEAN combined support ({{WEA_PIGEAN_COMBINED_CARDIO}}; $N$ = {{WEA_PIGEAN_COMBINED_CARDIO_N:,}}), one of the few areas where adding direct genetic evidence reduced performance. This suggests that pathway biology, including ion channels, contractile machinery, and lipid metabolism, is a stronger predictor of cardiovascular drug success than direct GWAS association. In endocrine disorders, PIGEAN indirect support likewise led ($\text{EA}|\text{RS}{=}1$ = {{WEA_PIGEAN_INDIRECT_ENDO}}; $N$ = {{WEA_PIGEAN_INDIRECT_ENDO_N}}).

% P4c: Areas where OTP is competitive
OTP performed competitively in therapeutic areas with well-characterized druggable landscapes and deep literature coverage. In immunology, OTP indirect support achieved $\text{EA}|\text{RS}{=}1$ = {{WEA_OTP_INDIRECT_IMMUNE}} ($N$ = {{WEA_OTP_INDIRECT_IMMUNE_N:,}}), far exceeding PIGEAN indirect support ({{WEA_PIGEAN_INDIRECT_IMMUNE}}; $N$ = {{WEA_PIGEAN_INDIRECT_IMMUNE_N:,}}). In metabolic disease, OTP combined ($\text{EA}|\text{RS}{=}1$ = {{WEA_OTP_COMBINED_METABOLIC:,}}; $N$ = {{WEA_OTP_COMBINED_METABOLIC_N}}) was narrowly ahead of PIGEAN combined ({{WEA_PIGEAN_COMBINED_METABOLIC:,}}; $N$ = {{WEA_PIGEAN_COMBINED_METABOLIC_N:,}}), though the $\text{EA}|\text{RS}{=}\text{Self}$ metric reversed this: PIGEAN combined ({{WEA_SELF_PIGEAN_COMBINED_METABOLIC:,}}) substantially outperformed OTP combined ({{WEA_SELF_OTP_COMBINED_METABOLIC}}), indicating that OTP has a strong pool of metabolic targets but PIGEAN ranks them more effectively. OTP also led in dermatology ($\text{EA}|\text{RS}{=}1$ = {{WEA_OTP_INDIRECT_DERM}}; $N$ = {{WEA_OTP_INDIRECT_DERM_N:,}}) and musculoskeletal disease ({{WEA_OTP_COMBINED_MUSCULO}}; $N$ = {{WEA_OTP_COMBINED_MUSCULO_N}}). In infection, all methods achieved negative or near-zero weighted EA (Figure~\ref{fig:weighted-ea}), consistent with the expectation that gene prioritization approaches are less applicable to infectious disease, where drug targets are often pathogen-derived rather than human genes.

% DISCUSSION: The oncology, neurology, cardiovascular, immune, and metabolic stories
% contain deeper mechanistic analysis that belongs in the Discussion section.


\input{sections/indirect-support/03302026/validation_clinical_03302026}

\subsection{A Phenome-Wide Map of Gene-Trait Relevance}

% STRUCTURE: This section summarizes the full phenome-wide application as a resource.
% It answers: what does the field get from PIGEAN that it didn't have before?

Applying the Mouse+MSigDB model across {{NUM_TOTAL_TRAITS_MOUSE_MSIGDB:,}} traits ({{NUM_COMMON_TRAITS_MOUSE_MSIGDB:,}} common, {{NUM_RARE_TRAITS_MOUSE_MSIGDB:,}} rare) and {{NUM_TOTAL_GENE_ANNOTATIONS:,}} gene annotations, PIGEAN produces a dense probabilistic map of gene-trait relevance. At a posterior probability threshold of 90\%, we identify {{NUM_GENETRAIT_REL_PP90:,}} gene-trait relationships, an average of {{NUM_PER_TRAIT_PP90}} genes per trait and {{NUM_PER_GENE_PP90}} traits per gene. At a more inclusive threshold of 50\%, this expands to {{NUM_GENETRAIT_REL_PP50:,}} relationships ({{NUM_PER_TRAIT_PP50}} per trait, {{NUM_PER_GENE_PP50}} per gene).

% TODO: Add comparison to existing resources - how many gene-trait relationships
% does Open Targets Genetics have? GWAS Catalog? Orphanet? This would quantify
% the "orders of magnitude denser" claim from the abstract.

Each gene-trait relationship is accompanied by provenance identifying the specific SNP associations or gene records contributing to direct support and the specific genesets contributing to indirect support. Across all traits, PIGEAN evaluates {{NUM_GENESET_TRAIT_ASSOC_MILLIONS}} million geneset-trait associations, retaining only those that survive the spike-and-slab regularization as non-zero contributors to indirect support.

% TODO: Add a brief description of what users can do with this resource.
% - Query a gene: which traits is it relevant to, and through which annotations?
% - Query a trait: which genes are relevant, and which are novel (indirect-only)?
% - Query a geneset: for which traits does it carry non-zero effect?
% - Bring your own annotations: run PIGEAN with a new CRISPR screen or single-cell atlas

% FIGURE SUGGESTION: A summary figure showing (a) density of gene-trait map vs. existing
% resources, (b) distribution of direct vs indirect contributions across gene-trait pairs,
% (c) example queries for a well-known gene (e.g., PCSK9) and a well-known trait (e.g., T2D).

% ANALYSIS IDEA: Show that the resource captures gene-trait relationships not in any
% existing curated database. How many PP>90% relationships have no corresponding
% entry in Open Targets or GWAS Catalog?

% STRUCTURE: This section presents two anecdotes of how the resource can generate mechanistic
% hypotheses: (1) the HLA-DRB1/MS sex paradox case, and (2) cross-referencing PIGEAN gene set
% enrichments with an independent CRISPR screen in human adipocytes.

% STRUCTURE: First anecdote -- querying PIGEAN geneset enrichments for a known gene-trait
% pair to generate mechanistic hypotheses that connect siloed literatures.

As an illustration of how the resource can generate mechanistic hypotheses, we examined the PIGEAN geneset enrichments for \textit{HLA-DRB1} in multiple sclerosis (MS). A collaborator had observed that male MS patients carrying \textit{HLA-DRB1} risk variants exhibited worse inflammation biomarkers (serum GFAP and NfL) and worse clinical outcomes than female carriers of the same variants---a pattern consistent with the well-documented MS sex paradox in which women develop MS approximately three times more often but men progress faster once diagnosed \cite{voskuhl_sex_2020}. Querying the PIGEAN resource for the top geneset enrichments of \textit{HLA-DRB1} in MS across two model configurations (a smaller MSigDB-only model and a larger Mouse+MSigDB model) revealed convergent biology: cell adhesion molecules (KEGG CAMs, enrichment score 0.848 in the small model), lymphotoxin-alpha (LTA, 0.562), common variable immunodeficiency (0.571), CD40 signaling (0.150), TNF receptor superfamily member~7/CD27 (0.204), and multiple gene ontology terms related to leukocyte adhesion and T~cell activation.

% TODO: MS_HLA_DRB1_NUM_GENESETS_SMALL - Number of non-zero genesets in the small model run
% TODO: MS_HLA_DRB1_NUM_GENESETS_LARGE - Number of non-zero genesets in the large model run
% TODO: Add formal references for each key paper cited in this section to references.bib

These enrichments decompose the \textit{HLA-DRB1} effect into a three-stage mechanistic cascade supported by independent experimental evidence: (1)~enhanced MHC-II-mediated antigen presentation driving T~cell priming, consistent with reports that \textit{HLA-DRB1*15:01} carriers show constitutively elevated HLA-DR expression on cortical microglia and larger grey matter lesions \cite{enz_increased_2020, yates_association_2015}; (2)~TNF superfamily amplification via LTA, CD40, and CD27 signaling that promotes B~cell activation and meningeal tertiary lymphoid tissue (TLT) formation, as demonstrated experimentally by Bates et al.\ (2022) and Naouar et al.\ (2026) \cite{bates_lymphotoxin_2022, naouar_lymphotoxin_2026}; and (3)~cell adhesion molecule-mediated lymphocyte transmigration across the blood-brain barrier, the rate-limiting step for CNS immune infiltration and the therapeutic target of natalizumab. Critically, each stage of this cascade interacts with sex-specific biology: estrogen partially suppresses TNF/LTA production by activated lymphocytes, providing upstream braking in women that is absent in men; blood-brain barrier endothelial responses to inflammation are sexually dimorphic \cite{weber_cerebrovascular_2021}; and downstream neuroprotection through estrogen receptor signaling on astrocytes and oligodendrocytes buffers women against the cortical damage driven by meningeal TLTs \cite{spence_neuroprotection_2011}. Men carrying \textit{HLA-DRB1} risk variants thus face a ``double hit'': the immune cascade runs with less hormonal restraint while the CNS neuroprotective buffer is weaker, consistent with the recent finding that \textit{HLA-DRB1*15:01} is associated with reduced longevity specifically in males \cite{mendoza_mejia_hla_2025}. While individual components of this model have been characterized in separate literatures---sex differences in MS progression, HLA-DRB1 effects on cortical pathology, and LTA-driven TLT formation---the integrated chain linking genotype to sex-differential severity through these specific intermediate mechanisms had not been previously proposed. The PIGEAN geneset enrichments independently recovered these nodes from GWAS data alone, illustrating how the resource can bridge siloed experimental literatures to generate testable mechanistic hypotheses.

% MISSING EVIDENCE: The sex-stratified analysis of HLA-DRB1 effects on GFAP/NfL
% would be strengthened by including the specific effect sizes from the collaborator's
% analysis once available. Consider adding variables:
% TODO: MS_HLA_DRB1_MALE_GFAP_BETA - Effect of DRB1 on GFAP in males
% TODO: MS_HLA_DRB1_FEMALE_GFAP_BETA - Effect of DRB1 on GFAP in females
% TODO: MS_HLA_DRB1_SEX_INTERACTION_P - P-value for sex interaction term

% ANALYSIS IDEA: Run PIGEAN on MS GWAS stratified by sex (if sex-stratified
% summary statistics are available) to test whether the geneset enrichment profile
% itself differs between male and female MS. This would directly test whether
% the model captures sex-differential biology at the geneset level.

% STRUCTURE: Second anecdote -- cross-referencing PIGEAN indirect support with an
% independent CRISPR screen to identify genes whose known function predicts a novel
% relationship with adipocyte biology.

\subsubsection{Cross-referencing indirect support with functional screens}

We next cross-referenced PIGEAN predictions with an independent high-content CRISPR screen in human adipocytes \cite{flannick_adipocyte_2024}. The screen profiled {{CRISPR_N_GENES}} genes across imaging phenotypes including lipid droplet morphology (BODIPY), actin cytoskeleton texture, and nuclear morphology (DAPI), clustering genes into nine factors. We queried the PIGEAN resource for these {{CRISPR_N_GENES}} genes across {{CRISPR_N_PHENOTYPES_MSIGDB}} curated metabolic, glycemic, hepatic, and anthropometric traits and classified each gene by its evidence type: \textit{understood mechanism} (high indirect and high direct support), \textit{novel GWAS discovery} (high direct, low indirect), or \textit{indirect-only} (high indirect support, low direct support). Of the screened genes, {{CRISPR_N_HITS}} ({{CRISPR_PCT_HITS}}\%) had at least one PIGEAN hit (combined support $\geq$ 2.0), including {{CRISPR_N_KNOWN}} of 42 known adipose-metabolism genes (\textit{PPARG} on rare diabetes, indirect support = 6.87; \textit{LPL} on TG:HDL ratio, indirect support = 6.61).

% TODO: CRISPR_N_GENES - Number of unique CRISPR screen genes (230)
% TODO: CRISPR_N_PHENOTYPES_MSIGDB - Number of metabolic phenotypes in msigdb model (537)
% TODO: CRISPR_N_HITS - Number of genes with at least one hit (114)
% TODO: CRISPR_PCT_HITS - Percentage of genes with hits (49.6)
% TODO: CRISPR_N_KNOWN - Number of known adipose genes with hits (40)
% TODO: CRISPR_N_NOVEL_INDIRECT - Number of novel indirect-only genes (5)

{{CRISPR_N_NOVEL_INDIRECT}} genes had high indirect support but low direct support for metabolic traits while also appearing in the CRISPR screen---genes whose pathway biology predicts metabolic relevance independently of any GWAS signal. \textit{EPS15} received the highest indirect support (3.64 for HbA1c adjusted for BMI; direct support = 1.30), driven by gene sets for decreased circulating insulin ($\beta$ = 1.17) and increased circulating glucose ($\beta$ = 1.02). EPS15 is an endocytic adaptor protein required for clathrin-coated pit assembly \cite{vanbergen_eps15_1995}, and the CRISPR screen assigned it to the lipid droplet morphology factor. In adipocytes, insulin-stimulated GLUT4 translocation to the plasma membrane depends on clathrin-dependent endocytic recycling \cite{leto_regulation_2012}; disruption of EPS15-dependent endocytosis would impair both insulin receptor internalization and GLUT4 re-exocytosis, consistent with the insulin/glucose gene set signal.

\textit{ADRA1B}, encoding the alpha-1B adrenergic receptor, received indirect support of 3.22 for fasting insulin adjusted for BMI (direct support = 1.11), driven by gene sets for abnormal glucose homeostasis ($\beta$ = 1.16), increased circulating insulin ($\beta$ = 1.10), and impaired glucose tolerance ($\beta$ = 1.07). The CRISPR screen placed it in the same lipid droplet morphology factor as \textit{EPS15}. Alpha-2 adrenergic receptors are established regulators of catecholamine-stimulated lipolysis in white adipose tissue \cite{lafontan_fat_2009, bartness_sympathetic_2014}, but a role for the alpha-1B subtype in adipocyte insulin sensitivity has not been reported. The convergence of a lipid droplet morphology phenotype in the screen with insulin/glucose gene set enrichments in PIGEAN suggests ADRA1B may modulate adipocyte insulin signaling through a mechanism distinct from classical lipolytic control.

\textit{RNF213}, an AAA+ ATPase and E3 ubiquitin ligase known as the major susceptibility gene for Moyamoya disease \cite{sugihara_rnf213_2019}, received indirect support of 2.30 for 2-hour glucose adjusted for BMI (direct support = 1.12), driven by gene sets for increased circulating glucose ($\beta$ = 1.32) and impaired glucose tolerance ($\beta$ = 1.12). The CRISPR screen assigned it to the actin/metabolic traits factor. Recent work has shown that RNF213 ubiquitinates proteins on the surface of intracellular lipid droplets \cite{ahel_rnf213_2023}, and the combination of this lipid droplet remodeling activity with PIGEAN-predicted glucose homeostasis relevance suggests RNF213 may regulate lipid droplet turnover in adipocytes through ubiquitin-dependent quality control. None of these three genes would have been prioritized by direct genetic association alone (all direct support $<$ 1.3).

% FIGURE SUGGESTION: A 2D scatter plot of indirect vs direct support for all 74 novel CRISPR
% genes, colored by evidence bucket (indirect-only, understood mechanism, novel GWAS,
% modest). Annotate the indirect-only genes. This would visually communicate the
% three-bucket framework and highlight the indirect-only quadrant.


\section{Discussion}\label{sec:discussion}

% This section currently contains mechanistic interpretation of the per-area
% weighted EA results from the drug target validation. Additional discussion
% paragraphs (e.g., model limitations, broader implications) can be added here.

\subsection{Therapeutic Area--Specific Performance of Indirect Support}

% DISCUSSION PARAGRAPH 1: Oncology -- literature saturation defeats OTP
% This paragraph explains WHY OTP anti-predicts oncology drug success.
The per-area analysis of weighted EA reveals that the relative performance of PIGEAN and literature-based resources varies systematically with the structure of each therapeutic area's druggable landscape. In oncology, where PIGEAN indirect support achieved the highest weighted EA of any area--method combination ($\text{EA}|\text{RS}{=}1$ = {{WEA_PIGEAN_INDIRECT_ONCO:,}}; $N$ = {{WEA_PIGEAN_INDIRECT_ONCO_N:,}}) and OTP combined was negative ({{WEA_OTP_COMBINED_ONCO:,}}; $N$ = {{WEA_OTP_COMBINED_ONCO_N:,}}; Figure~\ref{fig:weighted-ea}), the failure of literature-based scoring reflects a signal saturation problem. OTP assigns near-identical combined scores (0.95--0.999) to thousands of oncology gene--disease pairs because the cancer literature is vast: nearly every gene has been studied in the context of cancer, compressing OTP's score distribution and eliminating discriminative power. Consistent with this, OTP's top-half and bottom-half oncology targets have identical launched rates (7.1\% versus 7.1\%), meaning OTP's ranking provides zero predictive signal for which targets succeed clinically. PIGEAN's pathway-based indirect support, by contrast, scores genes based on their functional biology rather than publication frequency: PIGEAN's top-half oncology targets have a launched rate of 7.8\%, compared with 4.1\% for the bottom half, indicating genuine ranking discrimination. Among PIGEAN's top-ranked launched oncology targets are TP53 (head and neck carcinoma), SMO (basal cell carcinoma), JAK2 (polycythemia vera), MAP2K1 and MAP2K2 (thyroid carcinoma), BCL2 (chronic lymphocytic leukemia), and VEGFA (renal cell carcinoma), all targets identified through pathway membership rather than literature volume.

% DISCUSSION PARAGRAPH 2: Neurology -- cross-phenotype pathway integration
In neurology, where PIGEAN combined support achieved six-fold higher weighted EA than OTP ({{WEA_PIGEAN_COMBINED_NEURO:,}} versus {{WEA_OTP_COMBINED_NEURO}}; Figure~\ref{fig:weighted-ea}), the top gene sets driving indirect support reveal a coherent mechanistic basis for PIGEAN's advantage. The highest-contributing gene sets include amyloid-$\beta$ metabolic process (summed $\beta$ = 7.68), amyloid precursor protein metabolism (7.10), and $\alpha$-synuclein inclusion body (3.49), representing the core pathobiology of Alzheimer's and Parkinson's disease, alongside neuromuscular synapse morphology (7.64), ciliopathies (10.71), cilium assembly (6.71), mitochondrial DNA metabolic process (5.03), and Huntington's disease (2.07). This cross-phenotype integration is the distinguishing feature of indirect support: it borrows strength from related biological annotations, including neurodegenerative, neuromuscular, ciliary, and mitochondrial pathways, to prioritize targets for specific neurological indications. Literature co-occurrence, which scores genes independently for each disease term, cannot perform this kind of integrative reasoning across related phenotypes.

% DISCUSSION PARAGRAPH 3: Cardiovascular -- indirect support outperforms combined
The cardiovascular result, where PIGEAN indirect support (weighted EA = {{WEA_PIGEAN_INDIRECT_CARDIO}}; $N$ = {{WEA_PIGEAN_INDIRECT_CARDIO_N:,}}) exceeded combined support ({{WEA_PIGEAN_COMBINED_CARDIO}}; $N$ = {{WEA_PIGEAN_COMBINED_CARDIO_N:,}}; Figure~\ref{fig:weighted-ea}), provides direct evidence that pathway-based reasoning can outperform direct genetic association for drug target identification. The targets uniquely prioritized by indirect support, but not by combined support, are predominantly ion channel and calcium-handling genes with strong pathway biology but weak or absent GWAS signals: RYR2 for cardiac arrhythmias (indirect = 5.16, direct = 0.0), LCAT for myocardial infarction (indirect = 4.92, direct = 0.01), CACNA1C and CACNA1D for atrial fibrillation (indirect = 4.54 and 4.45; direct = 0.31 and 0.0), and KCNH2 for arrhythmias (indirect = 4.16, direct = 0.0), a launched antiarrhythmic target identified entirely by indirect evidence. The top gene sets for cardiovascular disease, including contractile fiber (summed $\beta$ = 7.83), arrhythmogenic right ventricular cardiomyopathy (7.75), cardiac conduction (5.56), cardiac fibrosis (5.68), cardiac hypertrophy (5.13), cholesterol metabolism (5.47), and plasma lipoprotein assembly and clearance (5.18), directly reflect the mechanistic pharmacology of cardiovascular drugs. Adding direct GWAS support dilutes this pathway signal by introducing genes with statistically significant but biologically ambiguous associations that do not map to established druggable biology.

% DISCUSSION PARAGRAPH 4: Immune -- OTP's literature depth captures a well-studied space
The immunology result presents an informative contrast: OTP indirect support achieved weighted EA = {{WEA_OTP_INDIRECT_IMMUNE}} ($N$ = {{WEA_OTP_INDIRECT_IMMUNE_N:,}}), while PIGEAN indirect support was near zero ({{WEA_PIGEAN_INDIRECT_IMMUNE}}; $N$ = {{WEA_PIGEAN_INDIRECT_IMMUNE_N:,}}; Figure~\ref{fig:weighted-ea}). OTP's top immune targets are the canonical cytokine and receptor drug targets: TNF (juvenile arthritis, launched), IL6 (rheumatoid arthritis, launched), IL17A (juvenile arthritis, launched), SERPING1 (hereditary angioedema, launched), and IGHE (asthma, Phase III), all thoroughly characterized in the immunology literature. PIGEAN's top immune targets, by contrast, are T-cell surface markers and signaling molecules (CD28, CD6, CD19, IL7R, CD200, LCK, CD3E) that are biologically central to immune function but whose clinical translation has been limited by safety and selectivity challenges. Only 6 of the top 50 genes ranked by PIGEAN and OTP overlapped. This pattern is consistent with the hypothesis that immunology drug development has been dominated by targeting well-characterized soluble mediators for which decades of literature faithfully record clinical translatability, whereas PIGEAN's pathway-based inference identifies central immune biology that is harder to exploit pharmacologically.

% DISCUSSION PARAGRAPH 5: Metabolic -- complementary ranking of lipid vs. glucose targets
In metabolic disease, OTP combined and PIGEAN combined achieved nearly identical weighted EA above RS = 1 ({{WEA_OTP_COMBINED_METABOLIC:,}} versus {{WEA_PIGEAN_COMBINED_METABOLIC:,}}; Figure~\ref{fig:weighted-ea}), but the $\text{EA}|\text{RS}{=}\text{Self}$ metric revealed that PIGEAN provides substantially better ranking within its target set ({{WEA_SELF_PIGEAN_COMBINED_METABOLIC:,}} versus {{WEA_SELF_OTP_COMBINED_METABOLIC}}). The ranking differences reflect the underlying evidence sources. PIGEAN assigns its highest scores to lipid metabolism targets: INS for type 2 diabetes (rank 1, combined = 12.3, launched), LDLR for hyperlipidemia (rank 2, combined = 11.8), PCSK9 for atherosclerosis (rank 69, combined = 10.0, launched), and NPC1L1 for hyperlipidemia (rank 77, combined = 9.87, launched), all targets supported by richly annotated cholesterol and lipoprotein pathway gene sets. OTP preferentially ranks glucose and incretin pathway targets: SLC5A2 for type 2 diabetes (OTP rank 14 versus PIGEAN rank 2,770, launched), GLP1R (OTP rank 49 versus PIGEAN rank 595, launched), and DPP4 (OTP rank 145 versus PIGEAN rank 1,810, launched). The discrepancy is consistent with the relative annotation depth of these pathways in MSigDB and the Mammalian Phenotype Ontology: lipid metabolism is extensively annotated across dozens of gene sets (cholesterol metabolism, lipoprotein pathways, sterol biosynthesis), whereas glucose transport and incretin signaling have sparser pathway representation. This finding suggests that the performance of indirect support is bounded by the completeness of the underlying gene set libraries, and that expanding annotation coverage of under-represented biological pathways could further improve PIGEAN's drug target prioritization.


\bibliography{references}% common bib file

% Methods section
\newpage
\input{sections/indirect-support/03302026/methods_03302026}

% Extended Data (figures and tables)
\newpage
\input{sections/indirect-support/03302026/extended_data_03302026}

% Supplementary Information (detailed derivations)
\newpage
\input{sections/indirect-support/03302026/supplementary_03302026}


\end{document}

% Some notes
% - pigean indirect + direct over common
% - pigean indirect over common (A2F and/or GWAS Cat)
% - combined
% - piccolo
% - OTG - All
% - magma 
% - bar or point
% - Temporal figure of indirect support regress magma2d
% - Take temporal direct support out 2c and cite in paper


% 3x2 for figure. Panel A is RS, Panel B is 1d

% We developed pigean and applied to XXXX common traits, predicted RS, and held when holding out future whatever. 
